“Ik breng rust in het veld!
Zo heeft men mij verteld.”
De kneus werd blootgesteld

na de vingerwijzing
aangaande de strekking
van Einfach’s deelneming.

De roedel grinnikte.
De vrager verklikte
wat het team hem flikte.

Hij begon te schreien.
Deze plagerijen
door de vogelvrijen

namen in omvang toe.
Men werd zijn zwakte moe.
Het bracht nood en gedoe.

Het nieuwe roedellid
zag hem als het doelwit
voor de herders exit

als de roedelleider.
Men zag hem als strijder
tegen de uitglijder

van de herder als broer.
Hij startte geen krachttoer,
maar reageerde stoer

door de opstandeling
te horen met achting.
Hij steunde een stemming

over Einfach’s vertrek.
Men verwenste de gek,
maar men hield toch zijn bek,

toen de nieuweling vroeg:
“Wie heeft van deze ploeg
van Einfach schoon genoeg?”

De indiener zijn haast
zat er mogelijk naast.
Hij opende verbaasd

dan maar zijn eigen mond,
die geen geluid uitzond.
Zijn strot bloedde verwond

door een mes van achter.
De stumper zijn wachter
bleek de vuige slachter.

Het verdriet ontwaakte,
toen ze een kreet slaakte,
omdat haar noot kraakte.

Haar knoop kon ontvouwen
en een knoop opbouwen
van lust en vertrouwen.

Haar vulva bleek tweeknoops
en werd zo iets goedkoops.
Emmy werd terloops loops.

De jager deed zetten.
Ze kon niet beletten
laat staan zich verzetten

tegen het ophitsen.
Haar brein liet zich gidsen.
Haar lijf ervoer flitsen

bij zijn verwoordingen
van veroveringen.
Zij beelde de dingen

in haar geestesoog uit
met gevoelens voluit,
via passend geluid.

Het verhaal prikkelde
en zijn stem tintelde.
Emmy bedisselde,

door een waanvoorstelling,
alsof de kieteling
niet leidde tot deining.

“Kathedralen van lichamen” is een pakkende en veelzeggende uitdrukking die zich leent voor diepgaande literaire en therapeutische reflectie. Binnen Emmy Nativus — een mystiek-erotische en autobiografische versroman — roept deze frase direct associaties op met het lichamelijke als heiligdom, als plek van aanbidding en transformatie.

**Analyse van de uitdrukking:**

1. **Logica en beeldspraak**
- De metafoor suggereert dat lichamen kathedralen zijn: groots, gewijd, gebouwd voor verering en introspectie.
- Kathedralen verwijzen naar architectuur, licht, resonantie, mysterie, en streven naar het hogere.
- Het lichamelijke wordt dus opgewaardeerd tot een tempel, een plaats van spiritueel of zintuiglijk beleven.

2. **Grammaticaal**
- “Kathedralen” is meervoud, wat duidt op collectiviteit — niet één lichaam, maar vele lichamen samen vormen een ensemble van kathedralen.
- “van lichamen” is een genitiefconstructie die eigendom, oorsprong of materiaal aangeeft.

3. **Historisch / Boekcorrect**
- In mystieke tradities wordt het lichaam vaak gezien als tempel of huis van de ziel (bijvoorbeeld Paulus’ brief: ‘Uw lichaam is een tempel van de heilige Geest’).
- Erotische spiritualiteit, zoals bij Georges Bataille en Teresa van Ávila, speelt met het idee van lichamelijke extase als poort naar het bovennatuurlijke.

4. **Stilistische en thematische suggesties**
- Binnen Emmy Nativus kan “kathedralen van lichamen” verwijzen naar scènes waarin fysieke ontmoeting en spirituele extase samenkomen.
- Het kan ook verwijzen naar het autobiografische: eigen lijf als literair bouwwerk waar herinnering, schaamte en verlangen samenkomen.

**Reflectievragen voor jouw creatieproces:**
- Hoe belichaamt de hoofdpersoon het idee van het lichaam als kathedraal?
- Zijn de lichamen in bijeenkomsten samenzang, harmonie, of eerder echo en eenzaamheid?
- Is er een contrast tussen de architectonische grootsheid van de kathedraal en de kwetsbaarheid van het lichamelijke?
- Welk soort muziek (melodie, zangstijl) zou passen bij een canto met deze titel?

**Rijm-, metrisch en thematisch suggesties** (niet uitgeschreven tenzij expliciet gevraagd):
- In welke canto of opera speelt deze metafoor een bepalende rol?
- Wil je het beeld uitwerken, muzikaal omlijsten, of alleen analyseren?

Laat gerust weten hoe je ‘kathedralen van lichamen’ verder wilt uitdiepen: wil je een songtekst, analyse, plotvoorstel, of een kant-en-klare canto uitwerking?

Verhalen van eendjes
met snavels als peentjes
en roeiende beentjes,

waar zij zonet waadde.
Alwaar hij zijn kwade
intenties verraadde,

nadat hij al zo vroeg
haar stiekem gadesloeg.
Hij zich daarbij afvroeg

wat hij kon casseren;
het zou opleveren
bij het inleveren

van Emmy bij de schout.
Haar vlieggewicht in goud.
Deze godsspraak was stout

vanuit de botterik.
Het bezorgde zo’n schrik
dat zij een ogenblik

zich verloor in de jacht.
De opgeilende kracht,
die gepaard ging met macht,

wilde zij gaan remmen.
Het wilde dier temmen
om het af te stemmen

op haar overwinning.
Het beest als verbeelding
gaf een hete spanning.

De in sex nog koter
werd zo geil als boter.
Het werd nog wat groter

met zijn vinger als steun.
Uit de wulpse miskleun
ontsnapte zacht gekreun

wat ze zwak intoomde,
hoewel ze sterk schroomde.
Haar geiligheid stroomde.

Ze wilde het geveol niet omarmen en kon ze niet overgeven eraan maar het lichaam overm,eesytede haar

Hij kon erom lachen
en ging door zigzaggen
over haar klit raggen.

Ze dicteerde zijn vijl
en ging weer voor de bijl.
Ze was onhoudbaar geil.

Ze voelde zich een hoer,
toen ze op de ijsvloer,
weer een climax ervoer.
Ze had een oneindig verlangen naar orgasmen die onbevredigbaar leken maar de jager wilde dat onderzoeken en bleef haar kietelen.

Toen woorden haar voerden
naar een grap, hoe woerden
soms een eend ontvoerden.

Dat die eerst niet wilde.
Van spanning eerst rilde,
daarna het uitgilde,

maar kalm bij de nasleep
zwom langs elke duckrape,
die zich aan haar vergreep.

“Is dat jouw hoop?”; dacht ze.
“Een onderdanige?
Speel ik die rol voor je!

Ben ik jouw proefkonijn.
Laat ik je in de schijn
dat ik bij je wil zijn.

Dat ik ervan genoot
dat ik helemaal bloot
me aan jouw blik aanbood.

Dat ik je lekker vind.
Dat voor mij als kleinkind
een opa me opwindt.

Daar ligt mijn verlossing.
Daarna ben jij naarling
slechts een herinnering.

Laat het animeren
dan snel arriveren.
De kou kan ik deren,

maar niet meer al te lang.”
De jager zijn aandrang
kwam zo in het gedrang

kon hij constateren.
De den met haar kleren
ging hij signaleren.

Met die blik kwam haar hoop
op een goede afloop.
De angst die haar bekroop

kwam voort uit de lusten
van hen die haar kusten
en het lieten rusten.

Of de oude jager
als hete belager
met een zoen als mager

aanbod genoegen nam,
was waarschijnlijk gezwam,
waar ze niet mee wegkwam.

Kon Emmy verwachten
haar leed te verzachten
als hij zou betrachten

wat terughoudendheid?
Zoals de jeugdigheid
deed uit onnozelheid?

Hij had haar aangeraakt,
haar onzuiver gemaakt
en haar gleufje gekraakt.

Volwassenen, aldus
Emmy haar oudste zus,
beginnen met een kus,

maar gaan dan snel verder.
Geen hoop op een redder.
Het werd dan precairder

en de curve steiler.
Hoe ouder hoe geiler.
Als dat gold als pijler

voor wat te wachten stond,
dan was haar zorg gegrond.
Emmy keek nogmaals rond.

Geen kans om te vluchten.
Ze had scha te duchten
wat betreft haar vruchten.

Ze keek naar de jager.
De oude belager
en grijze baarddrager

bleek een fitte opa.
Afzien en slecht karma,
lieten hun sporen na

in zijn slechte humeur
als rampspoed profiteur.
Haveloosheid als geur

ontwaarde ze langzaam
dampend uit zijn lichaam.
Wassen deed hij spaarzaam.

Er is wel erkenning
voor de bewondering
en de verwondering

waarmee hij haar voorzag
van een brede glimlach.
Om op zijn oude dag

nog getuige te zijn
van een meisje zo rein
en lieflijk in haar schijn.

Waar gaan we in dit boek
met Emmy naar op zoek?
Wat vormt in de driehoek

van lichaam, geest en ziel
in dit deel zijn profiel
onze achilleshiel?

De geest die onbeperkt
zijn duisternis doorwerkt?
En zo zijn ziel versterkt?

Dit werd vaak verboden
in die perioden,
die terreur echoden.

Bijvoorbeeld Frans Kafka
bedacht Gregor Samsa.
“Die Verwandlung” contra

Nazi-Duitsland kwam voort
omdat hij had verwoord
hoe macht geluk verstoord.

Hoe de menswaardigheid
en identiteit lijdt
onder autoriteit.

Of denk aan “Lolita”
met kindseks als thema.
Werd dit boek niet bijna

overal verboden.
Kunnen we die noden
door zwijgzaamheid doden?

Wordt iemand zijn natuur,
wel op de lange duur
verholpen door censuur?

Moest Lady Chatterley’s
intieme relatie
door seksificatie

worden uiteengezet?
Moest Lawrence schrijfset
door recht worden belet

zoals bij de Sade?
Wie helpt men dan daarmee?
Van wie komt het idee

dat voor monsters gelden
wat de schrijvers melden?
Of zijn ze juist helden

dat ze durven noemen
en zelfs in te zoomen
op wat we verbloemen?

Want had Dostojevski
met “Prestuplenie
i nakazanie” die

duisternis belicht?
De geest terecht beticht
van zijn onheilsgedicht?

De geest is vol met kwaad
en tot alles in staat
is wat Fjodor verraad.
